%=====================================================================
%  DAFTAR PUSTAKA
%  Menggunakan lingkungan thebibliography (tanpa BibTeX).
%  Urutan bibitem = urutan nomor sitasi pada dokumen.
%=====================================================================
\bagiantanpanomor{DAFTAR PUSTAKA}

\begingroup
\fontsize{8}{10}\selectfont
\begin{thebibliography}{99}

\bibitem{einstein1916}
A.~Einstein, ``Die Grundlage der allgemeinen Relativit\"{a}tstheorie,''
\textit{Annalen der Physik}, vol.~354, no.~7, hal.~769--822, 1916.

\bibitem{wald1984}
R.~M.~Wald, \textit{General Relativity}.
Chicago: University of Chicago Press, 1984.

\bibitem{schwarzschild1916}
K.~Schwarzschild, ``\"{U}ber das Gravitationsfeld eines Massenpunktes nach der
Einsteinschen Theorie,'' \textit{Sitzungsberichte der K\"{o}niglich Preussischen
Akademie der Wissenschaften zu Berlin}, hal.~189--196, 1916.

\bibitem{carroll2004}
S.~M.~Carroll, \textit{Spacetime and Geometry: An Introduction to General
Relativity}. San Francisco: Addison-Wesley, 2004.

\bibitem{frolov1998}
V.~P.~Frolov dan I.~D.~Novikov, \textit{Black Hole Physics: Basic Concepts and New
Developments}. Dordrecht: Springer, 1998.

\bibitem{misner1973}
C.~W.~Misner, K.~S.~Thorne, dan J.~A.~Wheeler, \textit{Gravitation}.
San Francisco: W.~H.~Freeman, 1973.

\bibitem{regge1957}
T.~Regge dan J.~A.~Wheeler, ``Stability of a Schwarzschild singularity,''
\textit{Physical Review}, vol.~108, no.~4, hal.~1063--1069, 1957.

\bibitem{vishveshwara1970}
C.~V.~Vishveshwara, ``Scattering of gravitational radiation by a Schwarzschild
black-hole,'' \textit{Nature}, vol.~227, hal.~936--938, 1970.

\bibitem{konoplya2011}
R.~A.~Konoplya dan A.~Zhidenko, ``Quasinormal modes of black holes: From
astrophysics to string theory,'' \textit{Reviews of Modern Physics}, vol.~83,
no.~3, hal.~793--836, 2011.

\bibitem{chandrasekhar1983}
S.~Chandrasekhar, \textit{The Mathematical Theory of Black Holes}.
Oxford: Oxford University Press, 1983.

\bibitem{berti2016}
E.~Berti, \textit{Black Hole Perturbation Theory}.
ICTP-SAIFR/IFT-UNESP Lecture Notes, 2016.

\bibitem{schutz2009}
B.~F.~Schutz, \textit{A First Course in General Relativity}, edisi ke-2.
Cambridge: Cambridge University Press, 2009.

\bibitem{kokkotas1999}
K.~D.~Kokkotas dan B.~G.~Schmidt, ``Quasinormal modes of stars and black holes,''
\textit{Living Reviews in Relativity}, vol.~2, no.~2, 1999.

\bibitem{nollert1999}
H.-P.~Nollert, ``Quasinormal modes: the characteristic `sound' of black holes and
neutron stars,'' \textit{Classical and Quantum Gravity}, vol.~16, no.~12,
hal.~R159--R216, 1999.

\bibitem{teukolsky1973}
S.~A.~Teukolsky, ``Perturbations of a rotating black hole. I. Fundamental equations
for gravitational, electromagnetic, and neutrino-field perturbations,''
\textit{The Astrophysical Journal}, vol.~185, hal.~635--647, 1973.

\bibitem{riley2006}
K.~F.~Riley, M.~P.~Hobson, dan S.~J.~Bence, \textit{Mathematical Methods for
Physics and Engineering}, edisi ke-3.
Cambridge: Cambridge University Press, 2006.

\end{thebibliography}
\endgroup
